\documentclass[11pt,a4paper]{article}
\usepackage[margin=1in]{geometry}
\usepackage{hyperref}
\usepackage{enumitem}
\usepackage{graphicx}
\usepackage{tikz}
\usepackage{amsmath}
\usepackage{amssymb}
\usepackage{booktabs}
\usepackage{titling}

\title{Alter}

\pretitle{\begin{center}\bfseries\fontsize{18bp}{18bp}\selectfont}
\posttitle{%
  \par\end{center}
  \begin{center}\large
  A Personal Digital Twin That Learns You\\[1.5em]
  Project Proposal for UCS503P\\
  Submitted to: Paramveer Sindhu\\
  \textbf{Thapar Institute of Engineering and Technology}
  \end{center}
  \vskip 0.5em%
}

\author{
  Lakkshya Jha\footnotemark[1]
  \and 
  Noorshan Singh\footnotemark[2]
  \and
  Saksham Raj\footnotemark[3]
}

\date{July 31, 2026}

\begin{document}
\maketitle
\footnotetext[1]{Roll No: 1024030708, Email: ljha_be24@thapar.edu}
\footnotetext[2]{Roll No: 1024030697, Email: nsingh_be24@thapar.edu}
\footnotetext[3]{Roll No: 1024030713, Email: sraj_be24@thapar.edu}

\section[Higher-order goal]{Higher-order goal \hfill \normalsize\normalfont\textit{...secondary framing}}
This project aims to transform human-computer interaction from generic assistance to deeply personalized, evolving digital companionship. While the topic of AI personalization is broad, the immediate engineering objective is to translate continuous learning and contextual memory retrieval into a measurable, deployable software solution.

\section[Time-to-value]{Time-to-value \hfill \normalsize\normalfont\textit{...primary heuristic}}
\begin{itemize}
    \item We will deliver meaningful increments quickly by prototyping the core memory retrieval workflow and validating it with users within a short iteration window.
    \item We will prioritize fast feedback over perfection by using weekly iterations (and targeting CI-backed automation early) to reduce release friction.
    \item Success is defined by what can be delivered and measured in the first iteration, then improved based on user feedback.
\end{itemize}

\section{Problem Statement}
Most AI assistants understand general information but have only shallow, fragmented understanding of the individual using them. Personal information is distributed across conversations, notes, documents, memories, decisions, habits, goals, and schedules. Common pain points include:
\begin{itemize}
    \item \textbf{Fragmentation:} Important experiences, decisions, and information are distributed across disconnected systems, causing lost context.
    \item \textbf{Low transparency:} Users cannot easily see or manage what the AI truly knows about them versus what it assumes.
    \item \textbf{Static Personalization:} Systems rely on fixed profiles or short-term context instead of continuously learning behavior.
\end{itemize}

\section[Proposed Solution]{Proposed Solution \hfill \normalsize\normalfont\textit{...Software Proposition}}

\subsection{Overview}
Build a web-based ``Personal Digital Twin'' system with the following components:
\begin{itemize}
    \item \textbf{User portal:} interface where users interact with Alter, chat, and submit daily notes or logs.
    \item \textbf{Memory layers:} Semantic, Episodic, and Behavioral memory modules.
    \item \textbf{Contextual retrieval engine:} automatically fetches relevant personal context for any query.
    \item \textbf{Prediction \& Adaptation engine:} simulates user decisions and adapts goals to user constraints.
    \item \textbf{Transparency dashboard:} allows users to view, edit, and control what Alter knows.
\end{itemize}

\subsection{Core workflow}
\begin{enumerate}
    \item User interacts with Alter through chat or notes.
    \item System categorizes and stores the interaction into Semantic, Episodic, or Behavioral memory.
    \item Upon a new query, the retrieval engine fetches relevant context from past data.
    \item The prediction engine generates a personalized response, simulating the user's likely decision or providing tailored advice.
\end{enumerate}

\subsection{Operational constraints for an educational, tier-2/3 setting}
\begin{itemize}
    \item Keep the solution usable on low-to-mid range devices via responsive web design.
    \item Avoid dependence on expensive hosted LLM APIs; focus on modular architecture and open-source models for NLP/vector retrieval where possible.
    \item Design for deployability using common web hosting and managed databases.
\end{itemize}

\section[Solution Approach]{Solution Approach \hfill \normalsize\normalfont\textit{...Engineering focus}}
This is an engineering course project, so we focus on scalable data processing, modular memory architecture, and retrieval workflows rather than creating novel foundation models.

\subsection[Duplicate/quality assistance]{Duplicate/quality assistance \hfill \normalsize\normalfont\textit{...non-research}}
Duplicate detection in memory storage will be heuristic-based:
\begin{itemize}
    \item Normalize text and categorize memories.
    \item Use a simple similarity measure (e.g., token overlap / TF-IDF cosine with threshold) without claiming state-of-the-art, to avoid storing duplicate facts or redundant episodic logs.
    \item Show top candidate duplicates to the user via the transparency dashboard for decision; never fully auto-delete or auto-reject.
\end{itemize}

\subsection[Web architecture]{Web architecture \hfill \normalsize\normalfont\textit{...web-based solution}}
A typical 3-tier web architecture:
\begin{itemize}
    \item \textbf{Frontend:} responsive UI in React/Next.js for interaction and memory management.
    \item \textbf{Backend API:} Python (FastAPI) for memory processing, retrieval engine, and interaction logic.
    \item \textbf{Database:} PostgreSQL with pgvector for storing user models and vector embeddings.
\end{itemize}

\subsection{CI/CD and fast delivery enabler}
CI/CD will be used to:
\begin{itemize}
    \item Automatically build and run tests on every change to the retrieval and memory components.
    \item Produce repeatable deployment artifacts to staging.
    \item Enable safe and frequent releases of incremental features.
\end{itemize}

\section[Evaluation Criterion]{Evaluation Criterion \hfill \normalsize\normalfont\textit{...measurable and attributable}}
We define success using measurable metrics tied to the core workflow.

\subsection{Primary evaluation metric}
\textbf{Context Retrieval Accuracy:} precision and recall of fetching the correct past memory for a given prompt.
\begin{itemize}
    \item Target: Top-K retrieval accuracy $\ge$ 85\% during pilot window.
    \item Attribution: measured using a test set of synthetic user interactions and queries.
\end{itemize}

\subsection{Secondary evaluation metrics}
\begin{itemize}
    \item \textbf{Prediction Accuracy:} percentage of times Alter accurately predicts user preferences or decisions based on historical data.
    \item \textbf{Response Time:} latency of the retrieval engine (target $\le$ 2 seconds).
    \item \textbf{User clarity:} user-reported understanding of why Alter made a specific prediction (using a simple 1--5 feedback question).
    \item \textbf{Reliability:} availability of staging/production for daily logging and querying ($\ge$ 99\% uptime).
\end{itemize}

\subsection[Pilot validation plan]{Pilot validation plan \hfill \normalsize\normalfont\textit{...high-level}}
\begin{itemize}
    \item Recruit 5--10 pilot users to interact with Alter daily over a 2-week period.
    \item Compare baseline (generic AI response) with Alter's contextual response.
    \item Log daily engagement and memory retrieval accuracy based on user feedback.
\end{itemize}

\section[Scalability]{Scalability \hfill \normalsize\normalfont\textit{...with foresight}}
The proposition should scale in both theoretical and practical senses.
\begin{enumerate}
    \item \textbf{Scalable (theoretically):} The system should support growth in the user's memory database by:
    \begin{itemize}
        \item decoupling components (frontend/backend/ML models),
        \item enabling efficient vector indexing (e.g., HNSW in pgvector) for fast similarity search,
        \item using stateless API design.
    \end{itemize}
    \item \textbf{Operationally deployable (practically):} The system should run on common web hosting:
    \begin{itemize}
        \item managed database,
        \item containerized or platform-hosted deployment,
        \item CI/CD pipeline for staging/production.
    \end{itemize}
\end{enumerate}

\section{Engine availability heuristic}
A mature set of ``engines'' (tooling/services) is available to implement this as a web-based project:
\begin{itemize}
    \item Web framework (React/Next.js) for REST APIs and templated/reactive pages.
    \item Managed relational or document database (PostgreSQL + pgvector).
    \item Authentication approach (token-based or session-based).
    \item NLP \& ML models (Sentence Transformers, spaCy) for text processing and embedding.
    \item Test framework for backend unit/integration tests.
    \item CI runner and staging deployment target.
\end{itemize}
We will use these mature components to focus on memory architecture, retrieval correctness, and personalized output rather than inventing new foundational models.

\section[Project scope and deliverables]{Project scope and deliverables \hfill \normalsize\normalfont\textit{...iteration-friendly}}

\subsection[Initial deliverable]{Initial deliverable \hfill \normalsize\normalfont\textit{...first iteration}}
\begin{itemize}
    \item User submission portal (chat/notes) + profile generation
    \item Basic Semantic and Episodic memory storage
    \item Contextual retrieval engine (vector search)
    \item Transparency dashboard: view stored memories
    \item CI: build + backend unit tests + linting
    \item CD to staging with a single command or pipeline job
\end{itemize}

\subsection{Subsequent deliverables}
\begin{itemize}
    \item Behavioral memory modeling and decision simulation
    \item Adaptive goal optimization and planning
    \item Automated duplicate/quality assistance in memory merging
    \item Improved search/filtering over personal timeline
\end{itemize}

\section{Risks and mitigations}
\begin{itemize}
    \item \textbf{Data privacy / consent:} minimize collected data; ensure users have full control to view, edit, and delete memories; use local or secure processing where possible.
    \item \textbf{User adoption:} provide an intuitive UI for daily logging; keep interaction friction minimal.
    \item \textbf{Evaluation measurement ambiguity:} clearly instrument accuracy metrics for memory retrieval using pre-labeled datasets before live user testing.
    \item \textbf{Operational issues in pilot:} use staging early; ensure CI/CD produces repeatable deployments.
\end{itemize}

\section{Summary}
This project proposes a deployable personal digital twin platform to transform generic AI into a continuously learning personalized assistant. It meets the course criteria by emphasizing time-to-value through rapid iterations, implementing CI/CD to support frequent safe releases, and using measurable evaluation criteria (especially context retrieval accuracy). It remains focused on engineering architecture, modular memory systems, and scalable vector retrieval rather than foundation model research.

\end{document}
